\section*{Chapitre \ref{ch:fm:hiring-and-compensation} --- Embauche et rémunération}

\begin{solution}{exo:fm:hiring-and-compensation:1}
$1-\Phi(1.2)=11.5~\%$ (environ un an sur neuf) et $\Phi(-1.0)=15.9~\%$ (un sur six).
\end{solution}

\begin{solution}{exo:fm:hiring-and-compensation:2}
$0.2\times(120-60-10)=10$ et $0.06\times120=7.2$.
\end{solution}

\begin{solution}{exo:fm:hiring-and-compensation:3}
Le cash de l'année, 0.6, plus un quart de la part différée de chacune des quatre attributions précédentes, $4\times0.1$~: une attribution entière de 1.
\end{solution}

\begin{solution}{exo:fm:hiring-and-compensation:4}
$s=4$~: part de chance $100/116=86.2~\%$, $T=100/(16\times0.5625)=11.1$ ans. $s=8$~: 61.0~\% et 2.78 ans.
\end{solution}

\begin{solution}{exo:fm:hiring-and-compensation:5}
Au-delà de 100~\% il faut l'approbation des actionnaires (jusqu'à 200~\%, à la majorité qualifiée). Au moins 40~\% doit être différé sur quatre à cinq ans (60~\% si
le montant est particulièrement élevé) et au moins 50~\% versé en actions ou en instruments équivalents.
\end{solution}

\begin{solution}{exo:fm:hiring-and-compensation:6}
La rémunération est $r\max(P\&L,0)$, convexe~: davantage de variance en relève l'espérance, alors que les pertes ne coûtent rien au trader sous zéro. Les limites plafonnent le risque~; le différé
assorti d'un \omterm{def:fm:hiring-and-compensation:malus}{malus} fait que les pertes ultérieures réduisent la rémunération antérieure.
\end{solution}

\begin{solution}{exo:fm:hiring-and-compensation:7}
Environ 0.76, contre 0.64 après cinq ans et 0.78 en forme fermée pour le P\&L moyen.
\end{solution}

\begin{solution}{exo:fm:hiring-and-compensation:8}
Une année à 25~millions de dollars est un faible indice de compétence (la chance représente 86~\% de la variance annuelle sur ce desk)~; une garantie paie quoi que donnent les années
suivantes, supprime le lien entre rémunération et résultats~; et là où s'appliquent les règles de l'UE, la rémunération variable garantie n'est permise qu'à l'embauche. Il faut proposer une \omterm{def:fm:hiring-and-compensation:deferral}{rémunération différée}
dont la valeur dépend des résultats futurs.
\end{solution}

\begin{solution}{pb:fm:hiring-and-compensation:1}
\begin{enumerate}
\item Voir la \cref{def:fm:hiring-and-compensation:structured}~; c'est la procédure de sélection la mieux classée dans la révision 2022 de la recherche sur la
validité prédictive.
\item Un jeu d'exercices chronométré, un entretien téléphonique structuré, une boucle d'\omterm{def:fm:hiring-and-compensation:structured}{entretiens structurés} sur site, notés indépendamment, et une réunion de décision sur
les notes~: les tests valides les moins coûteux d'abord.
\item La plupart des candidats ne seront pas embauchés, si bien qu'un test tardif consacre son effort à des personnes qu'un test précoce aurait pu écarter.
\item Voir la \cref{def:fm:hiring-and-compensation:pool}.
\item 6 pour un revenu de 100~; 0 pour un revenu de 70.
\item Elle ne verse rien une année où le desk ne dégage pas son coût du capital~; une enveloppe de revenus, elle, verse encore.
\item Voir la \cref{prop:fm:hiring-and-compensation:years}.
\item 86~\% et 11.1 ans.
\item 87~\%.
\item Une année~: 0.35, 0.34, 0.37~; cinq ans~: 0.64, 0.63, 0.61 (à la formule, ajusté du risque, discrétionnaire).
\item Le jugement porté sur le P\&L n'en est qu'une contraction et n'ajoute aucune information~; les observations sur la qualité des décisions, le risque et la contribution aux autres, elles, en apportent.
\item Environ 0.76.
\item Voir les \cref{def:fm:hiring-and-compensation:deferral,def:fm:hiring-and-compensation:malus,def:fm:hiring-and-compensation:buyout}.
\item Une attribution entière d'une année.
\item Le solde non acquis, soit une attribution annuelle, dans ses propres instruments différés, selon un calendrier similaire.
\item Rémunération variable plafonnée à 100~\% du fixe (200~\% avec approbation), au moins 40~\% différés sur quatre à cinq ans, au moins 50~\% en instruments~; le Royaume-Uni
a supprimé la limite de ratio à compter du 31 octobre 2023.
\item Le \omterm{def:fm:hiring-and-compensation:malus}{malus} réduit ce que l'entreprise détient encore~; la \omterm{def:fm:hiring-and-compensation:malus}{récupération} (clawback) doit recouvrer les sommes déjà versées.
\item Pas avec une garantie~: une année relève surtout de la chance~; il faut, le cas échéant, l'associer à une \omterm{def:fm:hiring-and-compensation:deferral}{rémunération différée} liée aux résultats futurs.
\item 87~\% de la variance de la \omterm{def:fm:hiring-and-compensation:pool}{rémunération à la formule} sur une année relève de la chance~; environ douze ans de P\&L (11.1 en forme fermée) sont nécessaires pour une corrélation de 0.8 avec la compétence.
\item Il faut verser une part du profit, après la \omterm{def:fm:limits-and-capital-allocation:raroc}{charge en capital}, répartie selon les résultats moyennés sur plusieurs années et la qualité observée des décisions, et différer une
grande partie assortie d'un \omterm{def:fm:hiring-and-compensation:malus}{malus}, de sorte que les résultats ultérieurs corrigent la rémunération antérieure.
\end{enumerate}
\end{solution}

\begin{solution}{iq:fm:hiring-and-compensation:1}
Exposer la rémunération aux pertes qui apparaissent plus tard (\omterm{def:fm:hiring-and-compensation:malus}{malus}), l'aligner sur des résultats de plus long terme et fidéliser.
\iqlookfor{\omterm{def:fm:hiring-and-compensation:malus}{malus} et fidélisation.}
\end{solution}

\begin{solution}{iq:fm:hiring-and-compensation:2}
Le \omterm{def:fm:hiring-and-compensation:malus}{malus} réduit les attributions non acquises~; la \omterm{def:fm:hiring-and-compensation:malus}{récupération} (clawback) reprend la rémunération déjà acquise ou versée.
\iqlookfor{la frontière de l'acquisition.}
\end{solution}

\begin{solution}{iq:fm:hiring-and-compensation:3}
Corrélation $\sqrt{0.16/(0.16+1/T)}$~; 0.8 exige $T=11.1$ ans.
\iqlookfor{rapport signal sur bruit et terme en $1/T$.}
\end{solution}

\begin{solution}{iq:fm:hiring-and-compensation:4}
Comparer la rémunération espérée sur plusieurs années, en incluant le solde non acquis perdu, les clauses d'acquisition et de \omterm{def:fm:hiring-and-compensation:malus}{malus}, ainsi que le capital et les limites qu'accorde chaque
entreprise.
\iqlookfor{le solde non acquis et les espérances ajustées du risque.}
\end{solution}

\begin{solution}{iq:fm:hiring-and-compensation:5}
La rémunération est convexe par rapport au P\&L~; on la maîtrise par des limites de risque, des \omterm{def:fm:limits-and-capital-allocation:raroc}{charges en capital} et un différé assorti d'un \omterm{def:fm:hiring-and-compensation:malus}{malus}.
\iqlookfor{la convexité et les remèdes.}
\end{solution}

\begin{solution}{iq:fm:hiring-and-compensation:6}
Ajouter des observations dont le rapport signal sur bruit est supérieur à celui du P\&L (qualité des décisions, calibration des prévisions, qualité de l'exécution), contracter le P\&L vers
elles, et différer pour que les résultats ultérieurs du P\&L puissent ajuster l'attribution.
\iqlookfor{davantage d'information, contraction et différé.}
\end{solution}
